\def\level{Terminale}  % Classe à droite
\def\course{NSI}
\def\chaptername{Programmation Orientée Objet} % Titre au centre
\def\docname{PR01}        % Identifiant à gauche

\pagestyle{cours_FP}
\makeheader{} % Affiche le bandeau


\begin{definition}[Programmation Orientée Objet - POO]{}
    La \textbf{Programmation Orientée Objet} ou \textbf{POO} est un paradigme de programmation informatique. Il consiste en la définition et l'interaction de briques logicielles appellées \verb+objets+. 

    Chaque objet est une \textbf{instance} de \textbf{classes} qui définissent le comportement et les propriétés des objets.
\end{definition}

\medskip

\begin{exemple}[Exemple de classes et d'objets]{\quad}{}
    \begin{itemize}
        \item \verb+Pokemon+ est une classe; c'est un concept. Pour le moment aucun pokemon n'existe mais juste l'idée;
        \item \verb+pikachu+ est un objet. C'est un pokemon particulier. On dit que c'est une instance de classe;
        \item une classe possède un nom et des caractéristiques (on parle d'\textbf{attributs});
        \item une classe possède des fonctionnalités (on parle de \textbf{méthodes})
        
    \end{itemize}
\end{exemple}

\begin{minipage}{0.45\linewidth}

    \begin{center}
    \includegraphics[width = 0.5\linewidth]{pikachu.png}
    \end{center}
    \begin{verbatim}
        Classe = Type d'objet
        Objet = Instance de classe
    \end{verbatim}
\end{minipage}
\begin{minipage}{0.54\linewidth}
    %\includegraphics[width = 1\linewidth]{drawit-diagram-37.png}
\begin{verbatim}
    Nom de la classe: Pokemon
    *******************************
    Attributs: 
        - nom (str)
        - points de vie: PV (int)
        - points de combat: PC (int)
        - attaque (str)
        - type (str)
    *******************************
    Méthodes:
        - connaitre les PV du pokemon
        - connaitre les PC du pokemon
        - modifier les PV du pokemon
        - savoir si le pokemon est toujours vivant
\end{verbatim}
\end{minipage}

\begin{propriete}[Diagramme de classe UML]{}

    \begin{minipage}{0.30\linewidth}
        \includegraphics[width = 0.90\linewidth]{poo01}

    \end{minipage}
    \begin{minipage}{0.20\linewidth}
        \includegraphics[width = 0.78\linewidth]{poo02}
    \end{minipage}\hfill
    \begin{minipage}{0.49\linewidth}
        \begin{itemize}
            \item \verb+Pokemon+ est le nom de la classe
            \item \verb+Pikachu+ est une instance de la classe \verb+Pokemon+
            \item \verb+pikachu+ est la valeur de l'attribut \verb+nom+ de l'instance \verb+Pikachu+
            \item \verb+60+ est la valeur de l'attribut \verb+pv+ de l'instance \verb+Pikachu+
        \end{itemize}
    
    \end{minipage}
    \end{propriete}

    \begin{propriete}{En python tout est objet}{}
    En python, toutes les variables sont déjà des objets. Pour vous en convaincre: \verb+dir(list)+

    Dans la liste des méthodes associées à la classe \verb+list+, on retrouve les méthodes connues (\verb+append+, \verb+pop+...) mais également des méthodes spéciales (ou \textbf{méthodes magiques}) qui sont entourées de double underscores (\verb+__init__+, \verb+__len__+, \verb+__add__+, etc.). Ces méthodes permettent de définir le comportement des objets de la classe.
\end{propriete}


\section{Construire une classe}

\begin{definition}[Construire une classe]{}
    Une classe est au minimum constituée d'un constructeur de classe.

    \begin{CodePythontexAlt}[Largeur=1\linewidth,Centre,Lignes=true,PremLigne=1]{}
        class Point:                            # Définition du nom de la classe
            def __init__(self, abs, ord):       # Constructeur de classe
                self.x = abs                    # Attribut x
                self.y = ord                    # Attribut y
    \end{CodePythontexAlt}
\end{definition}

    \begin{itemize}
        \item \verb+class Point+: c'est la porte d'entrée de la classe;
        \item \verb+def __init__(self, abs, ord)+: c'est la méthode qui est appelée au moment de la création d'une instance de la classe \verb+Point+;
        \item \verb+abs+ et \verb+ord+ sont des paramètre à entrer au moment de créer l'instance;
        \item \verb+self+: c'est l'objet en lui même. Pour le moment, il n'a pas de nom puisqu'on crée la classe. Ensuite \verb+self+ sera remplacé par le nom de l'instance de la classe;
        \item \verb+self.x+: on crée un attribut de classe \verb+x+ qui prend l'abscisse pour valeur;
        \item \verb+self.y+: on crée un attribut de classe \verb+y+ qui prend l'ordonnée pour valeur.

    \end{itemize}

    
\begin{definition}[Se servir d'une classe]{}
    
    Pour se servir d'une classe, il suffit maintenant de créer un objet (une instance) de la classe \verb+Point+

    \begin{CodePythontexAlt}[Largeur=1\linewidth,Centre,Lignes=true,PremLigne=1]{}
    A = Point(1, 7)     # A est une instance de la classe Point
    B = Point(0, 4)     # B est également une instance de la classe Point
\end{CodePythontexAlt}

Au moment de la création de l'instance A, le constructeur est exécuté. Dans toute la suite, \verb+self+ sera donc remplacé par \verb+A+.


\begin{CodePythontexAlt}[Largeur=1\linewidth,Centre,Lignes=true,PremLigne=1]{}
    print(A.x)     # On récupère l'attribut x de l'instance A. Affichage 1
    print(A.y)     # On récupère l'attribut y de l'instance A. Affichage 7

    dx = B.x - A.x    # On calcule la différence d'abscisse entre les points A et B
    dy = B.y - A.y    # On calcule la différence d'ordonnée entre les points A et B

    coeff_directeur = dy/dx   # On calcule le coefficient directeur de la droite (AB)
\end{CodePythontexAlt}

\end{definition}


\begin{definition}[Accesseurs et mutateurs]{}
En général on s'interdit d'accéder aux attributs de l'instance de la classe directement. A la place on crée des méthodes permettant d'accéder à ces attributs:

\begin{itemize}
    \item \textbf{Accesseurs} (getters): permet de récupérer la valeur de l'attribut;
    \item \textbf{Mutateurs} (setters): permet de modifier la valeur de l'attribut.
\end{itemize}

\begin{CodePythontexAlt}[Largeur=1\linewidth,Centre,Lignes=true,PremLigne=1]{}
    class Point:                                # Définition du nom de la classe
        def __init__(self, abs, ord):       # Constructeur de classe
            self.x = abs                    # Attribut x
            self.y = ord                    # Attribut y
        def get_coord(self):
            ''' Renvoie les coordonnées du point sous forme de tuple '''
            return (self.x, self.y)
        def set_coord(self, coord:tuple):
            ''' Modifie les coordonnées du point '''
            self.x, self.y = coord

    \end{CodePythontexAlt}

\end{definition}

\begin{easybox}[Affichage]{}
    Python permet d'afficher de deux manières: soit via la fonctionnalité \verb+print+ soit directement dans la console. Derrière ces fonctionnalités, il se cache deux méthodes:
    \begin{itemize}
        \item Pour la fonction \verb+print+, c'est la méthode \verb+__str__+
        \item Pour l'affichage dans la console c'est la méthode \verb+__repr__+
    \end{itemize}

    Par défaut ces méthodes sont crées lors de la création de la classe \verb+Point+ mais ne donnent rien d'intéressant:
        \begin{CodePythontexAlt}[Largeur=1\linewidth,Centre,Lignes=true,PremLigne=1]{}
        >>> print(A)
        <__main__.Point object at 0x00000233824A4D08>   
        # A est un objet de type Point situé à l'adresse mémoire 0x00000233824A4D08                 
    \end{CodePythontexAlt}

        On va donc surcharger (modifier) ces deux méthodes

        \begin{CodePythontexAlt}[Largeur=1\linewidth,Centre,Lignes=true,PremLigne=1]{}
            def __str__(self):
                ''' Surcharge la fonction print. Affichage 'Coordonnées = (abs;ord)'''
                return f'Coordonnées = ({self.x};{self.y})'
            def __repr__(self)
                ''' Surcharge __repr__. Affichage (abs;ord)'''
                return f'({self.x};{self.y})
        \end{CodePythontexAlt}
        


\end{easybox}

\begin{definition}[D'autres méthodes intéressantes]{}
    Dans le même ordre d'idée, il est possible de redéfinir des méthodes comme:
    \begin{itemize}
        \item \verb+__eq__(self, objet)+ qui renvoie \verb+True+ si \verb+self == objet+;
        \item \verb+__gt__(self, objet)+ qui renvoie \verb+True+ si \verb+self > objet+;
        \item \verb+__lt__(self, objet)+ qui renvoie \verb+True+ si \verb+self < objet+;
        \item \verb+__add__(self, objet)+ qui renvoie \verb-self + objet-;
        \item \verb+__mul__(self, objet)+ qui renvoie \verb+self * objet+;
    \end{itemize}
\end{definition}

\begin{propriete}[Utilisation des méthodes]{}

    \begin{CodePythontexAlt}[Largeur=1\linewidth,Centre,Lignes=true,PremLigne=1]{}
        >>> A = Point(1, 6)
        >>> B = Point(2, 7)
        >>> C = Point(1, 6)
        >>> B.set_coord((4, 8))
        >>> B.get_coord()
        (4, 8)
        >>> print(B)
        Coordonnées = (4;8)
        >>> B
        (4;8)
        >>> A == B
        False
        >>> A == C
        True
    \end{CodePythontexAlt}
    
\end{propriete}
\end{document}
\section{Les variables en Python}

\begin{easybox}{Déclarer une variable en Python}

	Pour déclarer une variable en Python, il suffit de lui donner un nom et de lui affecter une valeur à l'aide du symbole $=$. Par exemple, pour déclarer une variable \texttt{x} qui contient la valeur 10, on écrit:

\begin{minipage}{0.48\linewidth}
	\begin{CodePythontexAlt}[Largeur=1\linewidth,Centre,Lignes=true,PremLigne=1, Lignes=false,]{}
    a = 10
    \end{CodePythontexAlt}
\end{minipage}\hfill{}
\begin{minipage}{0.48\linewidth}
\begin{scratch}
\blocklook{Affecter à \selectmenu{x} la valeur \ovalnum{10}}\end{scratch}
\end{minipage}
\end{easybox}


\begin{easybox}{Le nom des variables en Python}
	\begin{itemize}
		\item Elles doivent commencer par une lettre ou un underscore (\texttt{\_}).
		\item Elles peuvent contenir des lettres, des chiffres et des underscores.
		\item Elles ne peuvent pas être des mots réservés du langage Python (comme \texttt{if}, \texttt{for}, \texttt{while}, etc.).
		\item \texttt{Variable} et \texttt{variable} sont considérées comme deux variables différentes.
	\end{itemize}
\end{easybox}

\begin{exemple}[Noms de variables valides et invalides]{}

	  \renewcommand{\arraystretch}{1.}
    \begin{tabular}{|l|p{12cm}|}
    \hline
    Bonne variable & \verb+ma_variable+ \\ \hline
    Mauvaise variable & \verb+Ma_variable+,  \verb+ma variable+, \verb+1_variable+, \verb+ variable@+ \\ \hline
    A éviter    & \verb+Variable non explicites+ \\ \hline
    \end{tabular}
\end{exemple}

\begin{easybox}{Afficher une variable en Python}

	Pour afficher la valeur d'une variable en Python, on utilise la fonction \texttt{print()}. Par exemple, pour afficher la valeur de la variable \texttt{x}, on écrit:

\begin{minipage}{0.48\linewidth}
	\begin{CodePythontexAlt}[Largeur=1\linewidth,Centre,Lignes=true,PremLigne=1, Lignes=false,]{}
    print(x)
    \end{CodePythontexAlt}
\end{minipage}\hfill{}
\begin{minipage}{0.48\linewidth}
\begin{scratch}
\blocklook{Afficher la valeur de \selectmenu{x}}\end{scratch}
\end{minipage}
\end{easybox}

\begin{easybox}{Les types de variables en Python}
	\begin{itemize}
		\item \textbf{int}: pour les nombres entiers (ex: 1, 42, -5)
		\item \textbf{float}: pour les nombres à virgule flottante (ex: 3.14, -0.001)
		\item \textbf{str}: pour les chaînes de caractères (ex: "Bonjour", 'Python')
		\item \textbf{bool}: pour les valeurs booléennes (ex: True, False)
	\end{itemize}
\end{easybox}

\begin{minipage}{0.19\linewidth}
\begin{lstlisting}[numbers=none]
# type (int)
a = 5
b = 6
\end{lstlisting}
    \end{minipage}\hfill
    \begin{minipage}{0.19\linewidth}    
\begin{lstlisting}[numbers=none]   
# type (str)
a = 'Bonjour'
b = '6'   
\end{lstlisting}    
    \end{minipage}\hfill
    \begin{minipage}{0.19\linewidth}    
\begin{lstlisting}[numbers=none] 
# type (float)
a = 5.0 
b = 6.8
\end{lstlisting}    
    \end{minipage}\hfill{}
\begin{minipage}{0.19\linewidth}    
\begin{lstlisting}[numbers=none]  
# type (bool)
a = True
b = False
\end{lstlisting}    
    \end{minipage}

\begin{easybox}{Opérations mathématiques de base en Python}{}
	\begin{minipage}{0.43\linewidth}
		\begin{itemize}[]
        \item[] \texttt{+} \quad  addition:  \verb~5 + 3~ donne 8
        \item[] \texttt{-} \quad  soustraction \verb~5 - 3~ donne 2
        \item[] \texttt{*} \quad  multiplication \verb~5 * 3~ donne 15
        \item[] \texttt{/} \quad  division \verb~6 / 2~ donne 3.0 (\texttt{float})
\end{itemize}
	\end{minipage}\hfill{}
	\begin{minipage}{0.55\linewidth}
\begin{itemize}[]
        \item[] \texttt{**} \quad  puissance \verb~2 ** 3~ donne 8
        \item[] \texttt{\%} \quad  reste de la division entière \texttt{7 \% 3} donne 1
        \item[] \texttt{//} \quad  quotient de la division entière \texttt{7 // 3} donne 2
    \end{itemize}
	\end{minipage}
\end{easybox}

\begin{minipage}{0.3\linewidth}
    \begin{ConsolePythontex}[Label=false]{}
        print(17+2)
        print(17-2)
    \end{ConsolePythontex}
\end{minipage}\hfill
\begin{minipage}{0.3\linewidth}
    \begin{ConsolePythontex}[Label=false]{}
		print(17*2)
		print(17/2)
    \end{ConsolePythontex}
\end{minipage}\hfill
\begin{minipage}{0.3\linewidth}
    \begin{ConsolePythontex}[Label=false]{}
        print(17%2)
        print(17//2)
    \end{ConsolePythontex}
\end{minipage}

\begin{exercice_cours}[Repérer les les types de variables et leurs valeurs dans les programmes suivants]
	\begin{minipage}{0.48\linewidth}
		\begin{itemize}
		\item \verb~x = 5~
		\item \verb~y = "Hello"~
		\item \verb~z = True~
		\item \verb~a = 3.14~
	\end{itemize}
	\end{minipage}\hfill{}
	\begin{minipage}{0.48\linewidth}
	\begin{itemize}
		\item \verb~c = 6/2~
		\item \verb~d = 6//2~
		\item \verb~f = '1' + '2'~
		\item \verb~g = 2**3~
	\end{itemize}
\end{minipage}
	
\end{exercice_cours}


\medskip

\section{Les variables de type chaine de caractère \texttt{str}}


\begin{easybox}{Les chaines de caractères}{}
    \smallskip
    Les chaines de caractères sont des variables qui contiennent du texte. Elles sont délimitées par des guillemets simples ou doubles.

    \begin{minipage}{0.32\linewidth}
        \begin{ConsolePythontex}[ Label=false]{}
            a = 'Bonjour'
            b = ' a tous'
            c = a + b
            print(c)
        \end{ConsolePythontex}
    \end{minipage}\hfill
    \begin{minipage}{0.32\linewidth}
        \begin{ConsolePythontex}[ Label=false]{}
            a = 'Bonjour'
            b = ' a tous'
            c = a * 3
            print(c)
        \end{ConsolePythontex}
    \end{minipage}\hfill
    \begin{minipage}{0.32\linewidth}
        \begin{ConsolePythontex}[ Label=false]{}
            a = 'Bonjour'
            b = ' a tous'
            c = a[0]
            print(c)
        \end{ConsolePythontex}
    \end{minipage}

    \medskip

    Dans l'exemple ci-dessus, la variable \verb+c+ contient la concaténation de \verb+a+ et \verb+b+ dans le premier cas, la répétition de \verb+a+ trois fois dans le deuxième cas et le premier caractère de \verb+a+ dans le troisième cas.
\end{easybox}

\begin{easybox}{indexation des chaines de caractères}{}
    
    \begin{minipage}{0.5\linewidth}
 Les chaines de caractères peuvent être vues comme un ensemble de caractères ``collés`` les uns aux autres

    Chaque caractère d'une chaine de caractères est indexé: 
    il y a un index qui donne sa position dans la chaine       
    \end{minipage}\hfill{}
    \begin{minipage}{0.5\linewidth}
 \begin{center}
        \begin{tabular}{|c|c|c|c|c|c|c|c|}
        \hline
        \textbf{Indice} & 0 & 1 & 2 & 3 & 4 & 5 & 6 \\
        \hline
        \textbf{Caractère} & B & o & n & j & o & u & r \\
        \hline
        \textbf{Indice Négatif} & -7 & -6 & -5 & -4 & -3 & -2 & -1 \\
        \hline
        \end{tabular}
    \end{center}       
    \end{minipage}
    
   
    

    \begin{minipage}[t]{0.48\linewidth}
        \begin{ConsolePythontex}[ Label=false]{}
            mot = 'Bonjour'
            mot[0]
            mot[4]
        \end{ConsolePythontex}
    \end{minipage}\hfill
    \begin{minipage}[t]{0.48\linewidth}
        \begin{ConsolePythontex}[ Label=false]{}
mot = 'Bonjour'
n = len(mot)
n
mot[-1]
		\end{ConsolePythontex}
    \end{minipage}
   
\end{easybox}

\begin{easybox}{Les opérations sur les chaines de caractères}{}   
    \smallskip
    \begin{itemize}
        \item[] \texttt{+} \quad pour la concaténation:  \verb~'Bonjour' + ' a tous'~ donne \verb~'Bonjour a tous'~
        \item[] \texttt{*} \quad pour la répétition:  \verb~'Bonjour' * 3~ donne \verb~'BonjourBonjourBonjour'~
        \item[] \texttt{[]} \quad pour l'indexation:  \verb~'Bonjour'[0]~ donne \verb~'B'~
        \item[] \texttt{in} \quad pour le test d'appartenance:  \verb~'on' in 'Bonjour'~ donne \verb~True~
        \item[] \texttt{not in} \quad pour le test d'appartenance:  \verb~'on' not in 'Bonjour'~ donne \verb~False~
        \item[] \texttt{len()} \quad pour la longueur:  \verb~len('Bonjour')~ donne \verb~7~
       
    \end{itemize}
\end{easybox}


\begin{minipage}[t]{0.32\linewidth}
    \begin{ConsolePythontex}[ Label=false]{}
        mot1 = 'Bonjour'
        mot2 = 'Thomas'
        mot1 + ' ' + mot2
    \end{ConsolePythontex}
\end{minipage}\hfill
\begin{minipage}[t]{0.32\linewidth}
    \begin{ConsolePythontex}[ Label=false]{}
        'j' in 'Bonjour'
        'b' in 'Bonjour'
    \end{ConsolePythontex}
\end{minipage}\hfill{}
\begin{minipage}[t]{0.32\linewidth}
    \begin{ConsolePythontex}[ Label=false]{}
        'abracadabra'[4]
        len('abracadabra')
    \end{ConsolePythontex}
\end{minipage}

\begin{easybox}{Opérations avancées sur les chaines de caractères: le Slicing}{}
\bigskip
    \renewcommand{\arraystretch}{1.2}
    \begin{tabular}{p{2cm}|p{13cm}}
        \texttt{[a:b]} & conserve les caractères d'indice \texttt{[a:b]}:  \verb~'Bonjour'[1:4]~ donne \verb~'onj'~ \\
        \texttt{[a:]} & conserve les caractères d'indice \texttt{[a:]}:  \verb~'Bonjour'[3:]~ donne \verb~'jour'~ \\
        \texttt{[:b]} & conserve les caractères d'indice \texttt{[:b]}:  \verb~'Bonjour'[:3]~ donne \verb~'Bon'~ \\
        \texttt{[::-1]} & inverse les caractères:  \verb~'Bonjour'[::-1]~ donne \verb~'ruojnoB'~ \\
        \texttt{[1:-1]} & élimine le premier et dernier caractère:  \verb~'Bonjour'[1:-1]~ donne \verb~'onjou'~ \\
    \end{tabular}
\end{easybox}

\begin{exercice_cours}[Prédire ce que vont afficher les codes suivants]{}

    \begin{minipage}[t]{0.32\linewidth}        
\begin{lstlisting}[numbers=none]
>>> mot = 'Numerique'
>>> n = len(mot)
>>> mot[5]
>>> mot[n]
>>> mot[-1]
>>> mot[8]
\end{lstlisting}
    \end{minipage} \hfill
    \begin{minipage}[t]{0.32\linewidth}        
\begin{lstlisting}[numbers=none]
>>> mot = 'Numerique'
>>> m in mot
>>> e not in mot
>>> mot+mot[2]
>>> mot*2
>>> 'abc'+mot
\end{lstlisting}
    \end{minipage}\hfill
\begin{minipage}[t]{0.32\linewidth}        
\begin{lstlisting}[numbers=none] 
>>> mot = 'Numerique'
>>> mot[1:5]
>>> mot[:4]
>>> mot[3:]
>>> mot[1:-1]
\end{lstlisting}
    \end{minipage}
\end{exercice_cours}

\begin{easybox}{Principe du transtypage}{}
    Les types des variables en python sont dynamiques, on peut modifier sa valeur ainsi que son type.
    
    \begin{minipage}[t]{0.48\linewidth}
        \begin{ConsolePythontex}[ Label=false]{}
            a = '5'     # a est de type str
            b = int(a)  # b est de type int
            print(b)
        \end{ConsolePythontex}
    \end{minipage}\hfill
    \begin{minipage}[t]{0.48\linewidth}
        \begin{ConsolePythontex}[ Label=false]{}
            a = 4       # a est de type int
            b = str(a)  # b est de type str
            print(b)
        \end{ConsolePythontex}
    \end{minipage}
    
    \end{easybox}

    \section{Les variables de type booléen \texttt{bool}}

    
\begin{easybox}{Les booléens}{}
    \smallskip
    Les booléens sont des variables qui ne peuvent prendre que deux valeurs: \verb+True+ (vrai) ou \verb+False+ (faux).
    
    \smallskip
    Les booléens peuvent être utilisés dans des opérations logiques.
    
	\medskip

    \begin{minipage}{0.32\linewidth}
        \begin{ConsolePythontex}[ Label=false]{}
a = True
b = False
c = a and b
c
        \end{ConsolePythontex}
    \end{minipage}\hfill
    \begin{minipage}{0.32\linewidth}
        \begin{ConsolePythontex}[ Label=false]{}
a = True
b = False
c = a or b
c
        \end{ConsolePythontex}
    \end{minipage}\hfill
    \begin{minipage}{0.32\linewidth}
        \begin{ConsolePythontex}[ Label=false]{}
a = True
c = not a
c
        \end{ConsolePythontex}
    \end{minipage}    
\end{easybox}

\medskip

\begin{minipage}[t]{6cm}
    \begin{center}Table de vérité du \verb+not+\end{center}
    \begin{longtable}[]{@{}cc@{}}
        \hline Expression & Résultat\tabularnewline\hline\endhead{}
        \texttt{not} True & False\tabularnewline{}
        \texttt{not} False & True\tabularnewline\hline{}
    \end{longtable}
\end{minipage}
\begin{minipage}[t]{6cm}
    \begin{center}Table de vérité du \verb+and+\end{center}
    \begin{longtable}[]{@{}cc@{}}
        \hline Expression & Résultat\tabularnewline\hline\endhead{}
        True \texttt{and} True & True\tabularnewline{}
        True \texttt{and} False & False\tabularnewline{}
        False \texttt{and} True & False\tabularnewline{}
        False \texttt{and} False & False\tabularnewline\hline{}
    \end{longtable}
\end{minipage}
\begin{minipage}[t]{6cm}
    \begin{center}Table de vérité du \verb+or+\end{center}
    \begin{longtable}[]{@{}cc@{}}
        \hline Expression & Résultat\tabularnewline\hline\endhead{}
        True \texttt{or} True & True\tabularnewline{}
        True \texttt{or} False & True\tabularnewline{}
        False \texttt{or} True & True\tabularnewline{}
        False \texttt{or} False & False\tabularnewline\hline{}
    \end{longtable}
\end{minipage}

\pagebreak

\begin{easybox}{Booléen comme résultat de comparaison}{}
    Les booléens sont le résultat d'une comparaison. Les opérateurs de comparaison sont:
    \smallskip
    \begin{itemize}
        \begin{minipage}{0.44\linewidth}
            \item[] \texttt{==} \quad pour le test d'égalité
            \item[] \texttt{!=} \quad pour le test de différence
            \item[] \texttt{<} \quad pour le test d'infériorité
        \end{minipage}
        \begin{minipage}{0.44\linewidth}
            \item[] \texttt{>} \quad pour le test de supériorité
            \item[] \texttt{<=} \quad pour le test d'infériorité ou d'égalité
            \item[] \texttt{>=} \quad pour le test de supériorité ou d'égalité
        \end{minipage}
    \end{itemize}

    \end{easybox}


\begin{easybox}{Booléen comme résultat de comparaison}{}  
    \smallskip
    \begin{minipage}{0.32\linewidth}    
        \begin{ConsolePythontex}[ Label=false]{}
a = 5
b = 6
c = (a == b)
c
\end{ConsolePythontex}
    \end{minipage}\hfill
    \begin{minipage}{0.32\linewidth}
        \begin{ConsolePythontex}[ Label=false]{}
a = 5
b = 6
c = (a != b)
c
        \end{ConsolePythontex}
    \end{minipage}\hfill
    \begin{minipage}{0.32\linewidth}
        \begin{ConsolePythontex}[ Label=false]{}
a = 5
b = 6
c = (a < b)
c
        \end{ConsolePythontex}
    \end{minipage}
    
\end{easybox}


\begin{exercice_cours}{    Quelle est la valeur de la variable \verb+c+ après l'exécution des codes suivants?}{}


    \begin{minipage}[t]{0.4\linewidth}
\begin{lstlisting}[numbers=none]   
>>> a = 3
>>> b = 1
>>> c = a > b
\end{lstlisting}

\begin{lstlisting}[numbers=none]
>>> a = True
>>> b = False
>>> c = (a or b) and (not a)
\end{lstlisting}

    \end{minipage} \hfill
    \begin{minipage}[t]{0.58\linewidth}
\begin{lstlisting}[numbers=none]
>>> a, b = True, False
>>> c = (a and b) or (not a) 
\end{lstlisting}   

\begin{lstlisting}[numbers=none]
>>> a = True
>>> b = False
>>> x, y = 7, 9
>>> c = (b or (x == y)) or (a and (x != y))
\end{lstlisting}  
    \end{minipage}

\end{exercice_cours}
